\section{齐：富庶东方与稷下的另一条国家道路}

从临淄向东，陆路通向半岛，水路和盐泽又把齐接到海边；向西，则是黄河下游与诸侯往来的道路。这样的地势不能自动换成胜兵，却能不断汇来粮食、货物与人。战国的齐因此不只是一支军队或一座都城：它有足以吸引游士的市场和俸禄，有可供征发的城邑，也有一旦失去盟友便格外显眼的财富。田氏取得齐国以后，正是在这块富庶的底盘上，把新的家族统治、列国外交和文化声望接到了一处；它后来遭遇的崩解，也从这里显出代价。

\begin{event}{公元前386年}{田氏获周册命，齐国易姓}{可靠纪年}{临淄与周王室}\eventanchor{WQ-0386-TIANQI}{战国·田氏代齐}
田氏取代姜氏，并非一日之间城门易帜。《史记·田敬仲完世家》把田和安置齐康公、向周廷求封与获命写在同一条权力转换的脉络中；具体步骤和先后未必能逐项复原，至少可以确定，到前386年，周王室承认田氏为齐侯，旧姜姓国君已不能再支配齐国的军政资源。临淄仍是那座临淄，城邑、田土与通往海滨的利路也没有随姓氏改变；改变的是谁能以齐侯之名分配这些资源。

这层名分仍值得争取。它和\eventref{WQ-0403-THREEJIN}{战国·三家分晋}所见的情形相似：周廷保有授名的仪式权威，却无力逆转列国中已经形成的强弱。田氏长期居于齐政的中心，积累的并不只是一个家族的声望，还有同卿大夫、城邑与民众相联的支持条件；周命为这一既成局面添上可以通行于列国之间的名号。新的统治者因而面对的并非如何从零建立国家，而是如何使接手的财富和人力不再只为田氏的胜利服务，而能成为齐国持久的安全。
\end{event}

齐在田氏统治下曾以富强自居，然而东方的富强也给邻国提供了共同防备它的理由。北方燕国在前四世纪末的内乱，使这种张力先以一次看似有利的干预显露出来。

\begin{event}{公元前314年}{齐军入燕，胜利留下北方的创伤}{可靠纪年}{燕地}\eventanchor{WQ-0314-QI-YAN}{战国·齐伐燕}
燕王哙、子之与太子一系的冲突撕裂了燕的命令链条，齐宣王遂遣兵北上。内乱中的道路、城邑和粮食供给本已失序，齐军得以深入，局势却没有因一支外军的到来而获得新的稳定。燕人抵抗，列国亦不愿看见齐把北方危机化作常久的领土；齐军终究撤离，留下的是被战事扰动的燕地和难以弥合的记忆。

《孟子·梁惠王下》保留了齐是否可以取燕的辩论，《史记·燕召公世家》则把它放进燕国由乱而复的叙事。两者成书背景不同，也都不能替我们精确还原一支支部队的行程；它们共同提示的，是这次干预的政治难题：以救乱或伐乱为名进入邻国，可以取得一时主动，却不能省去当地服从、列国承认和撤军之后的安排。前314年的损伤遂成为\eventref{WQ-YAN-ZHAO}{战国·燕昭王蓄力}得以凝聚人心的政治资源，也是后来\eventref{WQ-0284-QI}{战国·燕破齐}的近因之一。
\end{event}

燕昭王把雪耻、人材与对外联络慢慢编进国家准备，齐则在与诸国的竞逐中逐渐失去可用的余地。富庶能供给军队，也会使对手更愿意暂时搁置彼此的猜忌。

\begin{event}{公元前284年}{五国军渡济西，齐的防线崩解}{可靠纪年}{济西、临淄及齐地}\eventanchor{WQ-0284-QI}{战国·燕破齐}
乐毅统率燕、赵、秦、韩、魏五国之军攻齐。《史记·乐毅列传》所记济西之战，使齐军主力首先瓦解；盟军随后深入，临淄失守，齐湣王出奔，齐地大部落入燕军控制。传世叙事说乐毅下齐七十余城、改置燕郡，城数、设置的范围和各地降附的节奏都不宜据此写得过于齐整，但占领已从边境的战败延伸到齐国腹心，则无可回避。

这不是燕国单独击倒一个富国的故事。燕为盟主而乐毅为将，是多年筹备的结果；赵、秦、韩、魏同时出兵，才是使能条件。各国未必共享同一种报复，却能共同利用齐孤立的时刻。齐的粮食、城邑和市场没有消失，反而成为攻取后的目标；它缺少的是让这些财富在危机中转化为外援、守备与可靠指挥的外交和军事关系。田氏齐由此失去战前能够左右列国局势的主动权。
\end{event}

占领的难处很快显露。盟军各返本国之后，燕军要在广阔的齐地征收、守城并应付未降的据点；莒与即墨的存在，说明齐国的政治共同体虽被压碎，尚未消散。昭王死后，远在占领区的乐毅又失去原来那层君臣信用，燕惠王以骑劫代之，乐毅奔赵。关于田单如何施行反间、惠王如何受说的铺陈，主要见于《史记·田单列传》等后出的叙事，不能把其中每一个传递消息的动作当作可核的战场记录；易将、乐毅离去与齐军转守为攻之间的关联，却是局势逆转不能略去的人事转折。

\begin{event}{约公元前279年至前278年}{田单组织反攻，齐地渐次收复}{年代细节有争议}{即墨及齐地}\eventanchor{WQ-0279-TIANDAN}{战国·田单复齐}
田单所依托的不是一支凭空出现的救国军，而是即墨守军、尚能动员的齐地城邑，以及燕军后方命令已生裂缝的时机。《史记》把火牛夜出写成复国的明亮一刻；这个场景之所以流传，正在于它把夜袭、惊扰和军心变化压缩进可见的牲畜、火光与城门。然而它不能承担整个复国过程的解释。收复城邑需要粮食、向导、地方归附与连续的军事行动，具体年月和推进次序在文献中也有差异。

较稳妥的结论是，齐在前279年前后恢复了大部分失地，燕的占领体系随之退却。田单的胜利表明，攻城所得未必等于可以长期统治的土地；也表明残存据点一旦接上地方资源和对手的政治裂缝，仍能重新联结为国家。只是\eventref{WQ-0284-QI}{战国·燕破齐}造成的损耗不能由一次反攻抹平：齐复国之后仍有城邑与财富，却再难恢复战前那种持续向外施压的战略位置。
\end{event}

\subsection{稷下：财富如何变成名望，又为何不能代替安全}

齐的另一种力量并不在城下。后世称为稷下的空间，很可能依托临淄的城市、俸禄与政治中心而存在：游士可以到这里议论王道、法术、阴阳和治国，也可以借齐的名望使自己的学说传向其他国家。它不应被想成有统一课程、固定学籍的现代大学；更接近由君主和权贵供养、以列第与称号维系的知识网络。人才在此既是建议者，也是齐国向天下展示其富足、礼遇与开放程度的媒介。

\begin{textwindow}[《史记》笔下的稷下]
《史记·田敬仲完世家》说稷下诸士“皆赐列第为上大夫，不治而议论”。这句话出自汉代史家的国别叙事，足以说明后世记忆中的齐廷曾以俸禄、居处和议论来聚士；它却不足以替我们列出稷下的常设人数、课程和每日制度。
\end{textwindow}

孟子游齐的经历、荀子与齐的关联，以及邹衍、淳于髡等名字在传世文献中的聚集，使这种网络不只是一个地名。对齐廷而言，供养辩士的收益并非立刻化作军阵：他们带来政策语言、跨国声望与识别人才的渠道；对游士而言，临淄的富庶也提供了比单一宗族更宽的出路。可是，能使学说和人材流动的开放性，同样会把齐的政治选择暴露在列国注视之下。名士能帮助国家说明何为善政，不能替它守住济西的渡口，更不能自动弥合一次对外干预留下的敌意。

\begin{sourcenote}[材料的范围]
田氏易姓与齐、燕战争，可用《史记》诸世家、列传和《孟子》互相校读，并以《资治通鉴》的编年为线索；这些皆为后出传世材料，尤其战役细节和人物言辞须保留其叙事加工的可能。稷下的制度形态更难凭一条材料坐实，现代研究多从《史记》《荀子》《盐铁论》等的不同成书层次、临淄城市考古与战国游士流动来讨论。由此可以确认齐曾以资源吸纳知识与声望，却不能把“稷下学宫”写成制度完备、面貌不变的学院。
\end{sourcenote}

\begin{turningpoint}[制度得失：财富、外交与安全]
田氏以周命把既成的家族优势转为列国承认，齐的城邑、海滨之利和临淄网络又使它能够吸纳军队与游士。这样的富庶扩大了选择，却没有取消选择的代价：前314年入燕取得的是短期干预的机会，同时积累了燕国重建的动力；前284年的联盟进军证明，财富若不能接上稳固外交和危机中的指挥体系，便会成为他国共同争夺的对象；田单的收复则揭示，占领与复国都取决于地方社会是否仍可被动员。稷下提供的是一种文化与人才的国家能力，不能替代安全；安全一旦崩解，俸禄、辩论与市场也须重新依附于能够守住道路、城邑和盟约的政治秩序。
\end{turningpoint}
